\begin{abstract}\label{abstract}

\noindent \textbf{Introduction.}
This study systematically investigates the efficacy and safety of anti-angiogenic therapies in preventing \AMD, a leading cause of vision loss, focusing on their impact on visual outcomes and treatment burden. It aims to assess their impact on visual outcomes and treatment burden, providing a comprehensive overview of current therapeutic strategies and their effectiveness in AMD management.\par

\noindent \textbf{Methods.}
An extensive literature search across PubMed, ScienceDirect, and Google Scholar yielded ten studies meeting the inclusion criteria of randomized controlled trials and prospective studies, evaluating the efficacy and safety in AMD prevention. Data extraction and synthesis adhered to established systematic review protocols.\par

\noindent \textbf{Results.}
Analysis of the ten studies revealed that Treat-and-Extend (T\&E or TREX) regimens, which involve extending the interval between treatments based on the patient's response, offer efficacy comparable to traditional monthly dosing, significantly reducing treatment frequency. Innovative combination therapies demonstrated marked improvements in structural and visual outcomes of AMD. Safety profiles across various dosing schedules consistently indicated a low incidence of adverse effects, highlighting the potential of anti-angiogenic therapies to improve patient well-being alongside notable advancements in AMD management.\par

\noindent \textbf{Conclusion.}
This review synthesizes critical insights into effective and safe anti-angiogenic treatment options for AMD prevention. Treat-and-Extend regimens emerge as personalized, patient-centered approaches, ensuring efficacy with reduced intervention frequencies. The promising results from combination therapies indicate potential future directions for optimizing treatment efficacy while prioritizing patient safety and satisfaction.

\end{abstract}